% Options for packages loaded elsewhere
\PassOptionsToPackage{unicode}{hyperref}
\PassOptionsToPackage{hyphens}{url}
\documentclass[
]{article}
\usepackage{xcolor}
\usepackage[margin=1in]{geometry}
\usepackage{amsmath,amssymb}
\setcounter{secnumdepth}{-\maxdimen} % remove section numbering
\usepackage{iftex}
\ifPDFTeX
  \usepackage[T1]{fontenc}
  \usepackage[utf8]{inputenc}
  \usepackage{textcomp} % provide euro and other symbols
\else % if luatex or xetex
  \usepackage{unicode-math} % this also loads fontspec
  \defaultfontfeatures{Scale=MatchLowercase}
  \defaultfontfeatures[\rmfamily]{Ligatures=TeX,Scale=1}
\fi
\usepackage{lmodern}
\ifPDFTeX\else
  % xetex/luatex font selection
\fi
% Use upquote if available, for straight quotes in verbatim environments
\IfFileExists{upquote.sty}{\usepackage{upquote}}{}
\IfFileExists{microtype.sty}{% use microtype if available
  \usepackage[]{microtype}
  \UseMicrotypeSet[protrusion]{basicmath} % disable protrusion for tt fonts
}{}
\makeatletter
\@ifundefined{KOMAClassName}{% if non-KOMA class
  \IfFileExists{parskip.sty}{%
    \usepackage{parskip}
  }{% else
    \setlength{\parindent}{0pt}
    \setlength{\parskip}{6pt plus 2pt minus 1pt}}
}{% if KOMA class
  \KOMAoptions{parskip=half}}
\makeatother
\usepackage{color}
\usepackage{fancyvrb}
\newcommand{\VerbBar}{|}
\newcommand{\VERB}{\Verb[commandchars=\\\{\}]}
\DefineVerbatimEnvironment{Highlighting}{Verbatim}{commandchars=\\\{\}}
% Add ',fontsize=\small' for more characters per line
\usepackage{framed}
\definecolor{shadecolor}{RGB}{248,248,248}
\newenvironment{Shaded}{\begin{snugshade}}{\end{snugshade}}
\newcommand{\AlertTok}[1]{\textcolor[rgb]{0.94,0.16,0.16}{#1}}
\newcommand{\AnnotationTok}[1]{\textcolor[rgb]{0.56,0.35,0.01}{\textbf{\textit{#1}}}}
\newcommand{\AttributeTok}[1]{\textcolor[rgb]{0.13,0.29,0.53}{#1}}
\newcommand{\BaseNTok}[1]{\textcolor[rgb]{0.00,0.00,0.81}{#1}}
\newcommand{\BuiltInTok}[1]{#1}
\newcommand{\CharTok}[1]{\textcolor[rgb]{0.31,0.60,0.02}{#1}}
\newcommand{\CommentTok}[1]{\textcolor[rgb]{0.56,0.35,0.01}{\textit{#1}}}
\newcommand{\CommentVarTok}[1]{\textcolor[rgb]{0.56,0.35,0.01}{\textbf{\textit{#1}}}}
\newcommand{\ConstantTok}[1]{\textcolor[rgb]{0.56,0.35,0.01}{#1}}
\newcommand{\ControlFlowTok}[1]{\textcolor[rgb]{0.13,0.29,0.53}{\textbf{#1}}}
\newcommand{\DataTypeTok}[1]{\textcolor[rgb]{0.13,0.29,0.53}{#1}}
\newcommand{\DecValTok}[1]{\textcolor[rgb]{0.00,0.00,0.81}{#1}}
\newcommand{\DocumentationTok}[1]{\textcolor[rgb]{0.56,0.35,0.01}{\textbf{\textit{#1}}}}
\newcommand{\ErrorTok}[1]{\textcolor[rgb]{0.64,0.00,0.00}{\textbf{#1}}}
\newcommand{\ExtensionTok}[1]{#1}
\newcommand{\FloatTok}[1]{\textcolor[rgb]{0.00,0.00,0.81}{#1}}
\newcommand{\FunctionTok}[1]{\textcolor[rgb]{0.13,0.29,0.53}{\textbf{#1}}}
\newcommand{\ImportTok}[1]{#1}
\newcommand{\InformationTok}[1]{\textcolor[rgb]{0.56,0.35,0.01}{\textbf{\textit{#1}}}}
\newcommand{\KeywordTok}[1]{\textcolor[rgb]{0.13,0.29,0.53}{\textbf{#1}}}
\newcommand{\NormalTok}[1]{#1}
\newcommand{\OperatorTok}[1]{\textcolor[rgb]{0.81,0.36,0.00}{\textbf{#1}}}
\newcommand{\OtherTok}[1]{\textcolor[rgb]{0.56,0.35,0.01}{#1}}
\newcommand{\PreprocessorTok}[1]{\textcolor[rgb]{0.56,0.35,0.01}{\textit{#1}}}
\newcommand{\RegionMarkerTok}[1]{#1}
\newcommand{\SpecialCharTok}[1]{\textcolor[rgb]{0.81,0.36,0.00}{\textbf{#1}}}
\newcommand{\SpecialStringTok}[1]{\textcolor[rgb]{0.31,0.60,0.02}{#1}}
\newcommand{\StringTok}[1]{\textcolor[rgb]{0.31,0.60,0.02}{#1}}
\newcommand{\VariableTok}[1]{\textcolor[rgb]{0.00,0.00,0.00}{#1}}
\newcommand{\VerbatimStringTok}[1]{\textcolor[rgb]{0.31,0.60,0.02}{#1}}
\newcommand{\WarningTok}[1]{\textcolor[rgb]{0.56,0.35,0.01}{\textbf{\textit{#1}}}}
\usepackage{graphicx}
\makeatletter
\newsavebox\pandoc@box
\newcommand*\pandocbounded[1]{% scales image to fit in text height/width
  \sbox\pandoc@box{#1}%
  \Gscale@div\@tempa{\textheight}{\dimexpr\ht\pandoc@box+\dp\pandoc@box\relax}%
  \Gscale@div\@tempb{\linewidth}{\wd\pandoc@box}%
  \ifdim\@tempb\p@<\@tempa\p@\let\@tempa\@tempb\fi% select the smaller of both
  \ifdim\@tempa\p@<\p@\scalebox{\@tempa}{\usebox\pandoc@box}%
  \else\usebox{\pandoc@box}%
  \fi%
}
% Set default figure placement to htbp
\def\fps@figure{htbp}
\makeatother
\setlength{\emergencystretch}{3em} % prevent overfull lines
\providecommand{\tightlist}{%
  \setlength{\itemsep}{0pt}\setlength{\parskip}{0pt}}
\usepackage{booktabs}
\usepackage{longtable}
\usepackage{array}
\usepackage{multirow}
\usepackage{wrapfig}
\usepackage{float}
\usepackage{colortbl}
\usepackage{pdflscape}
\usepackage{tabu}
\usepackage{threeparttable}
\usepackage{threeparttablex}
\usepackage[normalem]{ulem}
\usepackage{makecell}
\usepackage{xcolor}
\usepackage{bookmark}
\IfFileExists{xurl.sty}{\usepackage{xurl}}{} % add URL line breaks if available
\urlstyle{same}
\hypersetup{
  pdftitle={Pasaporte\_Lab5: Naive Bayes},
  pdfauthor={Alejandro Antón, Juan Fernando, Fernando Mendoza},
  hidelinks,
  pdfcreator={LaTeX via pandoc}}

\title{Pasaporte\_Lab5: Naive Bayes}
\author{Alejandro Antón, Juan Fernando, Fernando Mendoza}
\date{2026-03-27}

\begin{document}
\maketitle

\section{Preparación de Librerías y
Datos}\label{preparaciuxf3n-de-libreruxedas-y-datos}

\begin{Shaded}
\begin{Highlighting}[]
\CommentTok{\# Librerías necesarias}
\FunctionTok{library}\NormalTok{(tidyverse)}
\FunctionTok{library}\NormalTok{(e1071)          }
\FunctionTok{library}\NormalTok{(caret)          }
\FunctionTok{library}\NormalTok{(ggplot2)}
\FunctionTok{library}\NormalTok{(knitr)}
\FunctionTok{library}\NormalTok{(kableExtra)}

\CommentTok{\# Cargar datos}
\FunctionTok{load}\NormalTok{(}\StringTok{"listings.RData"}\NormalTok{)}
\end{Highlighting}
\end{Shaded}

\section{Preprocesamiento de Datos}\label{preprocesamiento-de-datos}

\begin{Shaded}
\begin{Highlighting}[]
\CommentTok{\# Limpiar la columna price: remover \textquotesingle{}$\textquotesingle{} y \textquotesingle{},\textquotesingle{} y convertir a numérico}
\NormalTok{df\_limpio }\OtherTok{\textless{}{-}}\NormalTok{ listings }\SpecialCharTok{\%\textgreater{}\%}
  \FunctionTok{mutate}\NormalTok{(}\AttributeTok{price\_num =} \FunctionTok{as.numeric}\NormalTok{(}\FunctionTok{gsub}\NormalTok{(}\StringTok{"[$,]"}\NormalTok{, }\StringTok{""}\NormalTok{, price))) }\SpecialCharTok{\%\textgreater{}\%}
  \FunctionTok{filter}\NormalTok{(}\SpecialCharTok{!}\FunctionTok{is.na}\NormalTok{(price\_num), price\_num }\SpecialCharTok{\textgreater{}} \DecValTok{0}\NormalTok{, price\_num }\SpecialCharTok{\textless{}} \DecValTok{2000}\NormalTok{)}

\CommentTok{\# Calcular cuartiles para crear categorías balanceadas}
\NormalTok{q }\OtherTok{\textless{}{-}} \FunctionTok{quantile}\NormalTok{(df\_limpio}\SpecialCharTok{$}\NormalTok{price\_num, }\AttributeTok{probs =} \FunctionTok{c}\NormalTok{(}\FloatTok{0.25}\NormalTok{, }\FloatTok{0.50}\NormalTok{, }\FloatTok{0.75}\NormalTok{))}

\CommentTok{\# Crear variable categórica y discretizar predictoras para Naive Bayes}
\NormalTok{datos }\OtherTok{\textless{}{-}}\NormalTok{ df\_limpio }\SpecialCharTok{\%\textgreater{}\%}
  \FunctionTok{mutate}\NormalTok{(}
    \AttributeTok{price\_cat =} \FunctionTok{case\_when}\NormalTok{(}
\NormalTok{      price\_num }\SpecialCharTok{\textless{}=}\NormalTok{ q[}\DecValTok{1}\NormalTok{]                    }\SpecialCharTok{\textasciitilde{}} \StringTok{"Bajo"}\NormalTok{,}
\NormalTok{      price\_num }\SpecialCharTok{\textgreater{}}\NormalTok{ q[}\DecValTok{1}\NormalTok{] }\SpecialCharTok{\&}\NormalTok{ price\_num }\SpecialCharTok{\textless{}=}\NormalTok{ q[}\DecValTok{2}\NormalTok{] }\SpecialCharTok{\textasciitilde{}} \StringTok{"Medio{-}Bajo"}\NormalTok{,}
\NormalTok{      price\_num }\SpecialCharTok{\textgreater{}}\NormalTok{ q[}\DecValTok{2}\NormalTok{] }\SpecialCharTok{\&}\NormalTok{ price\_num }\SpecialCharTok{\textless{}=}\NormalTok{ q[}\DecValTok{3}\NormalTok{] }\SpecialCharTok{\textasciitilde{}} \StringTok{"Medio{-}Alto"}\NormalTok{,}
\NormalTok{      price\_num }\SpecialCharTok{\textgreater{}}\NormalTok{ q[}\DecValTok{3}\NormalTok{]                     }\SpecialCharTok{\textasciitilde{}} \StringTok{"Alto"}
\NormalTok{    ),}
    \AttributeTok{price\_cat =} \FunctionTok{factor}\NormalTok{(price\_cat, }\AttributeTok{levels =} \FunctionTok{c}\NormalTok{(}\StringTok{"Bajo"}\NormalTok{, }\StringTok{"Medio{-}Bajo"}\NormalTok{, }\StringTok{"Medio{-}Alto"}\NormalTok{, }\StringTok{"Alto"}\NormalTok{)),}
    
    \CommentTok{\# Asegurar tipos correctos y discretizar}
    \AttributeTok{accommodates =} \FunctionTok{as.numeric}\NormalTok{(accommodates),}
    \AttributeTok{bathrooms    =} \FunctionTok{as.numeric}\NormalTok{(bathrooms),}
    \AttributeTok{bedrooms     =} \FunctionTok{as.numeric}\NormalTok{(bedrooms),}
    \AttributeTok{beds         =} \FunctionTok{as.numeric}\NormalTok{(beds),}
    \AttributeTok{minimum\_nights =} \FunctionTok{as.numeric}\NormalTok{(minimum\_nights),}
    \AttributeTok{room\_type    =} \FunctionTok{as.factor}\NormalTok{(room\_type),}
    
    \AttributeTok{accommodates\_cat =} \FunctionTok{cut}\NormalTok{(accommodates, }\AttributeTok{breaks =} \FunctionTok{c}\NormalTok{(}\DecValTok{0}\NormalTok{, }\DecValTok{2}\NormalTok{, }\DecValTok{4}\NormalTok{, }\DecValTok{6}\NormalTok{, }\ConstantTok{Inf}\NormalTok{), }\AttributeTok{labels =} \FunctionTok{c}\NormalTok{(}\StringTok{"1{-}2"}\NormalTok{, }\StringTok{"3{-}4"}\NormalTok{, }\StringTok{"5{-}6"}\NormalTok{, }\StringTok{"7+"}\NormalTok{), }\AttributeTok{right =} \ConstantTok{TRUE}\NormalTok{),}
    \AttributeTok{bedrooms\_cat     =} \FunctionTok{cut}\NormalTok{(}\FunctionTok{replace\_na}\NormalTok{(bedrooms, }\DecValTok{0}\NormalTok{), }\AttributeTok{breaks =} \FunctionTok{c}\NormalTok{(}\SpecialCharTok{{-}}\ConstantTok{Inf}\NormalTok{, }\DecValTok{0}\NormalTok{, }\DecValTok{1}\NormalTok{, }\DecValTok{2}\NormalTok{, }\DecValTok{3}\NormalTok{, }\ConstantTok{Inf}\NormalTok{), }\AttributeTok{labels =} \FunctionTok{c}\NormalTok{(}\StringTok{"0"}\NormalTok{, }\StringTok{"1"}\NormalTok{, }\StringTok{"2"}\NormalTok{, }\StringTok{"3"}\NormalTok{, }\StringTok{"4+"}\NormalTok{)),}
    \AttributeTok{beds\_cat         =} \FunctionTok{cut}\NormalTok{(}\FunctionTok{replace\_na}\NormalTok{(beds, }\DecValTok{1}\NormalTok{), }\AttributeTok{breaks =} \FunctionTok{c}\NormalTok{(}\SpecialCharTok{{-}}\ConstantTok{Inf}\NormalTok{, }\DecValTok{1}\NormalTok{, }\DecValTok{2}\NormalTok{, }\DecValTok{3}\NormalTok{, }\ConstantTok{Inf}\NormalTok{), }\AttributeTok{labels =} \FunctionTok{c}\NormalTok{(}\StringTok{"1"}\NormalTok{, }\StringTok{"2"}\NormalTok{, }\StringTok{"3"}\NormalTok{, }\StringTok{"4+"}\NormalTok{)),}
    \AttributeTok{bathrooms\_cat    =} \FunctionTok{cut}\NormalTok{(}\FunctionTok{replace\_na}\NormalTok{(bathrooms, }\DecValTok{1}\NormalTok{), }\AttributeTok{breaks =} \FunctionTok{c}\NormalTok{(}\SpecialCharTok{{-}}\ConstantTok{Inf}\NormalTok{, }\DecValTok{1}\NormalTok{, }\FloatTok{1.5}\NormalTok{, }\DecValTok{2}\NormalTok{, }\ConstantTok{Inf}\NormalTok{), }\AttributeTok{labels =} \FunctionTok{c}\NormalTok{(}\StringTok{"1"}\NormalTok{, }\StringTok{"1.5"}\NormalTok{, }\StringTok{"2"}\NormalTok{, }\StringTok{"2+"}\NormalTok{)),}
    \AttributeTok{min\_nights\_cat   =} \FunctionTok{cut}\NormalTok{(minimum\_nights, }\AttributeTok{breaks =} \FunctionTok{c}\NormalTok{(}\DecValTok{0}\NormalTok{, }\DecValTok{1}\NormalTok{, }\DecValTok{3}\NormalTok{, }\DecValTok{7}\NormalTok{, }\ConstantTok{Inf}\NormalTok{), }\AttributeTok{labels =} \FunctionTok{c}\NormalTok{(}\StringTok{"1"}\NormalTok{, }\StringTok{"2{-}3"}\NormalTok{, }\StringTok{"4{-}7"}\NormalTok{, }\StringTok{"8+"}\NormalTok{), }\AttributeTok{right =} \ConstantTok{TRUE}\NormalTok{)}
\NormalTok{  ) }\SpecialCharTok{\%\textgreater{}\%}
  \FunctionTok{select}\NormalTok{(}
\NormalTok{    price\_num, price\_cat, room\_type, accommodates\_cat,}
\NormalTok{    bathrooms\_cat, bedrooms\_cat, beds\_cat,}
\NormalTok{    min\_nights\_cat}
\NormalTok{  ) }\SpecialCharTok{\%\textgreater{}\%}
  \FunctionTok{drop\_na}\NormalTok{()}
\end{Highlighting}
\end{Shaded}

Para implementar el algoritmo de Naive Bayes de manera efectiva,
preprocesamos los datos transformando las variables numéricas en
categorías o rangos. Además, la variable respuesta del precio se dividió
en cuatro segmentos equitativos utilizando cuartiles (Bajo, Medio-Bajo,
Medio-Alto, Alto) para garantizar que el algoritmo tenga suficientes
ejemplos de cada clase para aprender sin desbalancearse.

\begin{center}\rule{0.5\linewidth}{0.5pt}\end{center}

\section{División de Entrenamiento y
Prueba}\label{divisiuxf3n-de-entrenamiento-y-prueba}

\begin{Shaded}
\begin{Highlighting}[]
\FunctionTok{set.seed}\NormalTok{(}\DecValTok{123}\NormalTok{)}

\CommentTok{\# Índices de entrenamiento: 70\%}
\NormalTok{train\_index }\OtherTok{\textless{}{-}} \FunctionTok{createDataPartition}\NormalTok{(datos}\SpecialCharTok{$}\NormalTok{price\_cat, }\AttributeTok{p =} \FloatTok{0.70}\NormalTok{, }\AttributeTok{list =} \ConstantTok{FALSE}\NormalTok{)}

\NormalTok{train\_data }\OtherTok{\textless{}{-}}\NormalTok{ datos[train\_index, ]}
\NormalTok{test\_data  }\OtherTok{\textless{}{-}}\NormalTok{ datos[}\SpecialCharTok{{-}}\NormalTok{train\_index, ]}

\FunctionTok{cat}\NormalTok{(}\StringTok{"Tamaño del conjunto de entrenamiento:"}\NormalTok{, }\FunctionTok{nrow}\NormalTok{(train\_data), }\StringTok{"}\SpecialCharTok{\textbackslash{}n}\StringTok{"}\NormalTok{)}
\end{Highlighting}
\end{Shaded}

\begin{verbatim}
## Tamaño del conjunto de entrenamiento: 51920
\end{verbatim}

\begin{Shaded}
\begin{Highlighting}[]
\FunctionTok{cat}\NormalTok{(}\StringTok{"Tamaño del conjunto de prueba:       "}\NormalTok{, }\FunctionTok{nrow}\NormalTok{(test\_data), }\StringTok{"}\SpecialCharTok{\textbackslash{}n}\StringTok{"}\NormalTok{)}
\end{Highlighting}
\end{Shaded}

\begin{verbatim}
## Tamaño del conjunto de prueba:        22249
\end{verbatim}

Mantenemos la partición del 70\% para entrenamiento y 30\% para prueba
utilizando la misma semilla (123) de los laboratorios anteriores. Esto
es un paso crucial para asegurar que la comparación final entre todos
los algoritmos sea válida y justa.

\begin{center}\rule{0.5\linewidth}{0.5pt}\end{center}

\section{Actividades 1 y 2: Entrenamiento y Evaluación Numérica
(Regresión)}\label{actividades-1-y-2-entrenamiento-y-evaluaciuxf3n-numuxe9rica-regresiuxf3n}

\begin{Shaded}
\begin{Highlighting}[]
\CommentTok{\# Entrenamos Naive Bayes con corrección de Laplace para evitar probabilidades cero}
\NormalTok{modelo\_nb }\OtherTok{\textless{}{-}} \FunctionTok{naiveBayes}\NormalTok{(}
\NormalTok{  price\_cat }\SpecialCharTok{\textasciitilde{}}\NormalTok{ room\_type }\SpecialCharTok{+}\NormalTok{ accommodates\_cat }\SpecialCharTok{+}\NormalTok{ bathrooms\_cat }\SpecialCharTok{+}\NormalTok{ bedrooms\_cat }\SpecialCharTok{+}\NormalTok{ beds\_cat }\SpecialCharTok{+}\NormalTok{ min\_nights\_cat,}
  \AttributeTok{data   =}\NormalTok{ train\_data,}
  \AttributeTok{laplace =} \DecValTok{1} 
\NormalTok{)}

\CommentTok{\# 1. Calculamos el valor representativo promedio de cada categoría}
\NormalTok{medias\_precio }\OtherTok{\textless{}{-}}\NormalTok{ train\_data }\SpecialCharTok{\%\textgreater{}\%}
  \FunctionTok{group\_by}\NormalTok{(price\_cat) }\SpecialCharTok{\%\textgreater{}\%}
  \FunctionTok{summarise}\NormalTok{(}\AttributeTok{precio\_promedio =} \FunctionTok{mean}\NormalTok{(price\_num, }\AttributeTok{na.rm =} \ConstantTok{TRUE}\NormalTok{))}

\CommentTok{\# 2. Extraemos las probabilidades de cada categoría para el conjunto de prueba}
\NormalTok{pred\_probabilidades }\OtherTok{\textless{}{-}} \FunctionTok{predict}\NormalTok{(modelo\_nb, }\AttributeTok{newdata =}\NormalTok{ test\_data, }\AttributeTok{type =} \StringTok{"raw"}\NormalTok{)}

\CommentTok{\# 3. Calculamos el "Precio Predicho" numérico (Valor Esperado)}
\NormalTok{precio\_predicho\_nb }\OtherTok{\textless{}{-}}\NormalTok{ (pred\_probabilidades[, }\StringTok{"Bajo"}\NormalTok{] }\SpecialCharTok{*}\NormalTok{ medias\_precio}\SpecialCharTok{$}\NormalTok{precio\_promedio[}\DecValTok{1}\NormalTok{]) }\SpecialCharTok{+}
\NormalTok{                      (pred\_probabilidades[, }\StringTok{"Medio{-}Bajo"}\NormalTok{] }\SpecialCharTok{*}\NormalTok{ medias\_precio}\SpecialCharTok{$}\NormalTok{precio\_promedio[}\DecValTok{2}\NormalTok{]) }\SpecialCharTok{+}
\NormalTok{                      (pred\_probabilidades[, }\StringTok{"Medio{-}Alto"}\NormalTok{] }\SpecialCharTok{*}\NormalTok{ medias\_precio}\SpecialCharTok{$}\NormalTok{precio\_promedio[}\DecValTok{3}\NormalTok{]) }\SpecialCharTok{+}
\NormalTok{                      (pred\_probabilidades[, }\StringTok{"Alto"}\NormalTok{] }\SpecialCharTok{*}\NormalTok{ medias\_precio}\SpecialCharTok{$}\NormalTok{precio\_promedio[}\DecValTok{4}\NormalTok{])}

\CommentTok{\# 4. Evaluamos el error de esta predicción continua}
\NormalTok{rmse\_nb\_reg }\OtherTok{\textless{}{-}} \FunctionTok{sqrt}\NormalTok{(}\FunctionTok{mean}\NormalTok{((test\_data}\SpecialCharTok{$}\NormalTok{price\_num }\SpecialCharTok{{-}}\NormalTok{ precio\_predicho\_nb)}\SpecialCharTok{\^{}}\DecValTok{2}\NormalTok{, }\AttributeTok{na.rm =} \ConstantTok{TRUE}\NormalTok{))}
\NormalTok{mae\_nb\_reg }\OtherTok{\textless{}{-}} \FunctionTok{mean}\NormalTok{(}\FunctionTok{abs}\NormalTok{(test\_data}\SpecialCharTok{$}\NormalTok{price\_num }\SpecialCharTok{{-}}\NormalTok{ precio\_predicho\_nb), }\AttributeTok{na.rm =} \ConstantTok{TRUE}\NormalTok{)}

\FunctionTok{cat}\NormalTok{(}\StringTok{"Métricas del seudo{-}modelo de Regresión Naive Bayes:}\SpecialCharTok{\textbackslash{}n}\StringTok{"}\NormalTok{)}
\end{Highlighting}
\end{Shaded}

\begin{verbatim}
## Métricas del seudo-modelo de Regresión Naive Bayes:
\end{verbatim}

\begin{Shaded}
\begin{Highlighting}[]
\FunctionTok{cat}\NormalTok{(}\StringTok{"RMSE:"}\NormalTok{, }\FunctionTok{round}\NormalTok{(rmse\_nb\_reg, }\DecValTok{2}\NormalTok{), }\StringTok{"}\SpecialCharTok{\textbackslash{}n}\StringTok{"}\NormalTok{)}
\end{Highlighting}
\end{Shaded}

\begin{verbatim}
## RMSE: 220.28
\end{verbatim}

\begin{Shaded}
\begin{Highlighting}[]
\FunctionTok{cat}\NormalTok{(}\StringTok{"MAE:"}\NormalTok{, }\FunctionTok{round}\NormalTok{(mae\_nb\_reg, }\DecValTok{2}\NormalTok{), }\StringTok{"}\SpecialCharTok{\textbackslash{}n}\StringTok{"}\NormalTok{)}
\end{Highlighting}
\end{Shaded}

\begin{verbatim}
## MAE: 137.72
\end{verbatim}

El algoritmo de Naive Bayes es de naturaleza probabilística y no calcula
regresiones continuas directas. Para evaluar su capacidad predictiva
sobre el precio exacto (Actividades 1 y 2), multiplicamos la
probabilidad que el algoritmo asignó a cada categoría por el precio
promedio de dicha categoría, obteniendo un valor esperado continuo. Los
resultados muestran un Error Cuadrático Medio (RMSE) bastante elevado,
lo que indica que el algoritmo falla por un margen considerable al
intentar adivinar la tarifa exacta por noche. Esto ocurre porque Naive
Bayes asume que todas las características de la casa son independientes
entre sí, lo cual no aplica en la realidad inmobiliaria donde el tamaño
y las camas están directamente relacionados.

\begin{center}\rule{0.5\linewidth}{0.5pt}\end{center}

\section{Actividad 3: Comparación Histórica de Modelos de
Regresión}\label{actividad-3-comparaciuxf3n-histuxf3rica-de-modelos-de-regresiuxf3n}

\begin{Shaded}
\begin{Highlighting}[]
\CommentTok{\# Valores aproximados obtenidos en laboratorios anteriores bajo los mismos datos}
\NormalTok{resultados\_regresion }\OtherTok{\textless{}{-}} \FunctionTok{data.frame}\NormalTok{(}
  \AttributeTok{Algoritmo =} \FunctionTok{c}\NormalTok{(}\StringTok{"Regresión Lineal Múltiple"}\NormalTok{, }\StringTok{"Árbol de Decisión (Prof. 4)"}\NormalTok{, }\StringTok{"Naive Bayes (Probabilístico)"}\NormalTok{),}
  \AttributeTok{RMSE\_Aproximado =} \FunctionTok{c}\NormalTok{(}\StringTok{"Muy Alto (MSE 13M)"}\NormalTok{, }\StringTok{"250.45"}\NormalTok{, }\FunctionTok{round}\NormalTok{(rmse\_nb\_reg, }\DecValTok{2}\NormalTok{))}
\NormalTok{)}

\FunctionTok{kable}\NormalTok{(resultados\_regresion, }\AttributeTok{caption =} \StringTok{"Comparación Histórica de Modelos Numéricos"}\NormalTok{) }\SpecialCharTok{\%\textgreater{}\%}
  \FunctionTok{kable\_styling}\NormalTok{(}\AttributeTok{bootstrap\_options =} \FunctionTok{c}\NormalTok{(}\StringTok{"striped"}\NormalTok{, }\StringTok{"hover"}\NormalTok{), }\AttributeTok{full\_width =} \ConstantTok{FALSE}\NormalTok{)}
\end{Highlighting}
\end{Shaded}

\begin{longtable}[t]{ll}
\caption{\label{tab:comparacion_regresion}Comparación Histórica de Modelos Numéricos}\\
\toprule
Algoritmo & RMSE\_Aproximado\\
\midrule
Regresión Lineal Múltiple & Muy Alto (MSE 13M)\\
Árbol de Decisión (Prof. 4) & 250.45\\
Naive Bayes (Probabilístico) & 220.28\\
\bottomrule
\end{longtable}

Al contrastar Naive Bayes con los modelos de las semanas previas, el
Árbol de Decisión de profundidad 4 se mantiene como la herramienta más
eficiente para predecir el precio exacto. La Regresión Lineal forzaba
relaciones rectas inexistentes, y Naive Bayes falla por asumir
independencia total. El Árbol de Regresión, al crear reglas específicas
(ej. ``Casa entera con más de 3 camas''), logra ajustarse mejor a la
variabilidad del mercado.

\begin{center}\rule{0.5\linewidth}{0.5pt}\end{center}

\section{Actividad 4: Modelo de Clasificación
(Categorías)}\label{actividad-4-modelo-de-clasificaciuxf3n-categoruxedas}

\begin{Shaded}
\begin{Highlighting}[]
\CommentTok{\# Extraemos las probabilidades a priori del modelo ya entrenado}
\NormalTok{apriori\_df }\OtherTok{\textless{}{-}} \FunctionTok{data.frame}\NormalTok{(}
  \AttributeTok{Categoria =} \FunctionTok{names}\NormalTok{(modelo\_nb}\SpecialCharTok{$}\NormalTok{apriori),}
  \AttributeTok{Probabilidad =} \FunctionTok{as.numeric}\NormalTok{(}\FunctionTok{prop.table}\NormalTok{(modelo\_nb}\SpecialCharTok{$}\NormalTok{apriori))}
\NormalTok{)}

\FunctionTok{ggplot}\NormalTok{(apriori\_df, }\FunctionTok{aes}\NormalTok{(}\AttributeTok{x =}\NormalTok{ Categoria, }\AttributeTok{y =}\NormalTok{ Probabilidad, }\AttributeTok{fill =}\NormalTok{ Categoria)) }\SpecialCharTok{+}
  \FunctionTok{geom\_col}\NormalTok{(}\AttributeTok{show.legend =} \ConstantTok{FALSE}\NormalTok{) }\SpecialCharTok{+}
  \FunctionTok{geom\_text}\NormalTok{(}\FunctionTok{aes}\NormalTok{(}\AttributeTok{label =}\NormalTok{ scales}\SpecialCharTok{::}\FunctionTok{percent}\NormalTok{(Probabilidad, }\AttributeTok{accuracy =} \FloatTok{0.1}\NormalTok{)), }\AttributeTok{vjust =} \SpecialCharTok{{-}}\FloatTok{0.5}\NormalTok{, }\AttributeTok{size =} \DecValTok{4}\NormalTok{) }\SpecialCharTok{+}
  \FunctionTok{scale\_fill\_brewer}\NormalTok{(}\AttributeTok{palette =} \StringTok{"Blues"}\NormalTok{) }\SpecialCharTok{+}
  \FunctionTok{labs}\NormalTok{(}\AttributeTok{title =} \StringTok{"Probabilidades A Priori por Categoría de Precio"}\NormalTok{,}
       \AttributeTok{x =} \StringTok{"Categoría de Precio"}\NormalTok{, }\AttributeTok{y =} \StringTok{"Probabilidad"}\NormalTok{) }\SpecialCharTok{+}
  \FunctionTok{theme\_minimal}\NormalTok{(}\AttributeTok{base\_size =} \DecValTok{13}\NormalTok{) }\SpecialCharTok{+}
  \FunctionTok{ylim}\NormalTok{(}\DecValTok{0}\NormalTok{, }\FloatTok{0.35}\NormalTok{)}
\end{Highlighting}
\end{Shaded}

\begin{center}\includegraphics{pasaporte_5_files/figure-latex/apriori-1} \end{center}

Al cambiar el enfoque hacia la clasificación, el objetivo es predecir a
qué rango de mercado pertenece una propiedad. El modelo entrenado
calcula las probabilidades previas (mostradas en la gráfica) reflejando
un balance perfecto de aproximadamente 25\% por categoría. Esto evita
que el algoritmo se vuelva sesgado y simplemente intente predecir
siempre la categoría con más datos.

\begin{center}\rule{0.5\linewidth}{0.5pt}\end{center}

\section{Actividades 5 y 6: Predicción, Eficiencia y Matriz de
Confusión}\label{actividades-5-y-6-predicciuxf3n-eficiencia-y-matriz-de-confusiuxf3n}

\begin{Shaded}
\begin{Highlighting}[]
\CommentTok{\# Predicciones de clase sobre datos no vistos}
\NormalTok{pred\_clase }\OtherTok{\textless{}{-}} \FunctionTok{predict}\NormalTok{(modelo\_nb, }\AttributeTok{newdata =}\NormalTok{ test\_data, }\AttributeTok{type =} \StringTok{"class"}\NormalTok{)}

\CommentTok{\# Generación de la matriz de confusión}
\NormalTok{conf\_matrix }\OtherTok{\textless{}{-}} \FunctionTok{confusionMatrix}\NormalTok{(pred\_clase, test\_data}\SpecialCharTok{$}\NormalTok{price\_cat)}

\CommentTok{\# Convertir a data frame para visualización gráfica}
\NormalTok{conf\_df }\OtherTok{\textless{}{-}} \FunctionTok{as.data.frame}\NormalTok{(conf\_matrix}\SpecialCharTok{$}\NormalTok{table)}
\FunctionTok{colnames}\NormalTok{(conf\_df) }\OtherTok{\textless{}{-}} \FunctionTok{c}\NormalTok{(}\StringTok{"Predicho"}\NormalTok{, }\StringTok{"Real"}\NormalTok{, }\StringTok{"Frecuencia"}\NormalTok{)}

\FunctionTok{ggplot}\NormalTok{(conf\_df, }\FunctionTok{aes}\NormalTok{(}\AttributeTok{x =}\NormalTok{ Real, }\AttributeTok{y =}\NormalTok{ Predicho, }\AttributeTok{fill =}\NormalTok{ Frecuencia)) }\SpecialCharTok{+}
  \FunctionTok{geom\_tile}\NormalTok{(}\AttributeTok{color =} \StringTok{"white"}\NormalTok{) }\SpecialCharTok{+}
  \FunctionTok{geom\_text}\NormalTok{(}\FunctionTok{aes}\NormalTok{(}\AttributeTok{label =}\NormalTok{ Frecuencia), }\AttributeTok{size =} \DecValTok{5}\NormalTok{, }\AttributeTok{fontface =} \StringTok{"bold"}\NormalTok{) }\SpecialCharTok{+}
  \FunctionTok{scale\_fill\_gradient}\NormalTok{(}\AttributeTok{low =} \StringTok{"\#E3F2FD"}\NormalTok{, }\AttributeTok{high =} \StringTok{"\#1565C0"}\NormalTok{) }\SpecialCharTok{+}
  \FunctionTok{labs}\NormalTok{(}\AttributeTok{title =} \StringTok{"Matriz de Confusión {-} Naive Bayes"}\NormalTok{,}
       \AttributeTok{subtitle =} \FunctionTok{paste}\NormalTok{(}\StringTok{"Exactitud Global:"}\NormalTok{, }\FunctionTok{round}\NormalTok{(conf\_matrix}\SpecialCharTok{$}\NormalTok{overall[}\StringTok{"Accuracy"}\NormalTok{] }\SpecialCharTok{*} \DecValTok{100}\NormalTok{, }\DecValTok{1}\NormalTok{), }\StringTok{"\%"}\NormalTok{),}
       \AttributeTok{x =} \StringTok{"Categoría Real"}\NormalTok{, }\AttributeTok{y =} \StringTok{"Categoría Predicha"}\NormalTok{, }\AttributeTok{fill =} \StringTok{"Frecuencia"}\NormalTok{) }\SpecialCharTok{+}
  \FunctionTok{theme\_minimal}\NormalTok{(}\AttributeTok{base\_size =} \DecValTok{13}\NormalTok{)}
\end{Highlighting}
\end{Shaded}

\begin{center}\includegraphics{pasaporte_5_files/figure-latex/confusion-1} \end{center}

\begin{Shaded}
\begin{Highlighting}[]
\CommentTok{\# Métricas detalladas por clase}
\NormalTok{metricas\_clase }\OtherTok{\textless{}{-}} \FunctionTok{as.data.frame}\NormalTok{(conf\_matrix}\SpecialCharTok{$}\NormalTok{byClass)}
\NormalTok{metricas\_clase}\SpecialCharTok{$}\NormalTok{Clase }\OtherTok{\textless{}{-}} \FunctionTok{rownames}\NormalTok{(metricas\_clase)}

\NormalTok{metricas\_clase }\SpecialCharTok{\%\textgreater{}\%}
  \FunctionTok{select}\NormalTok{(Clase, Sensitivity, Specificity, }\StringTok{\textasciigrave{}}\AttributeTok{Pos Pred Value}\StringTok{\textasciigrave{}}\NormalTok{, F1) }\SpecialCharTok{\%\textgreater{}\%}
  \FunctionTok{mutate}\NormalTok{(}\FunctionTok{across}\NormalTok{(}\FunctionTok{where}\NormalTok{(is.numeric), }\SpecialCharTok{\textasciitilde{}} \FunctionTok{round}\NormalTok{(.x, }\DecValTok{3}\NormalTok{))) }\SpecialCharTok{\%\textgreater{}\%}
  \FunctionTok{kable}\NormalTok{(}\AttributeTok{col.names =} \FunctionTok{c}\NormalTok{(}\StringTok{"Clase"}\NormalTok{, }\StringTok{"Sensibilidad"}\NormalTok{, }\StringTok{"Especificidad"}\NormalTok{, }\StringTok{"Precisión (VPP)"}\NormalTok{, }\StringTok{"F1{-}Score"}\NormalTok{),}
        \AttributeTok{caption =} \StringTok{"Métricas de Eficiencia por Categoría de Precio"}\NormalTok{) }\SpecialCharTok{\%\textgreater{}\%}
  \FunctionTok{kable\_styling}\NormalTok{(}\AttributeTok{bootstrap\_options =} \FunctionTok{c}\NormalTok{(}\StringTok{"striped"}\NormalTok{, }\StringTok{"hover"}\NormalTok{), }\AttributeTok{full\_width =} \ConstantTok{FALSE}\NormalTok{)}
\end{Highlighting}
\end{Shaded}

\begin{longtable}[t]{llrrrr}
\caption{\label{tab:metricas_clase}Métricas de Eficiencia por Categoría de Precio}\\
\toprule
 & Clase & Sensibilidad & Especificidad & Precisión (VPP) & F1-Score\\
\midrule
Class: Bajo & Class: Bajo & 0.670 & 0.781 & 0.505 & 0.576\\
Class: Medio-Bajo & Class: Medio-Bajo & 0.328 & 0.822 & 0.381 & 0.352\\
Class: Medio-Alto & Class: Medio-Alto & 0.258 & 0.850 & 0.365 & 0.302\\
Class: Alto & Class: Alto & 0.628 & 0.841 & 0.567 & 0.596\\
\bottomrule
\end{longtable}

Al poner a prueba el modelo con datos nuevos, obtenemos una radiografía
completa de sus fortalezas y debilidades a través de la matriz de
confusión:

\begin{itemize}
\tightlist
\item
  \textbf{Dónde se equivocó menos (Efectividad):} El algoritmo tiene una
  gran capacidad para identificar los extremos. Rara vez confunde una
  propiedad económica (``Bajo'') con una propiedad premium (``Alto'').
  La sensibilidad es más fuerte en estos polos porque sus
  características estructurales son muy marcadas y distintas.
\item
  \textbf{Dónde se equivocó más (Confusión):} Los bloques centrales de
  la matriz evidencian que el modelo tiene serios problemas para
  diferenciar entre los rangos intermedios (``Medio-Bajo'' y
  ``Medio-Alto''). Al basarnos en cuartiles, la diferencia de precio
  entre estas propiedades puede ser de apenas un par de dólares, por lo
  que sus características físicas (baños, camas) se solapan fuertemente.
\item
  \textbf{Importancia de los errores:} Para SmartStay, sugerir un precio
  Medio-Bajo a una propiedad que realmente es Medio-Alta representa una
  pérdida de ingresos moderada pero manejable. Sin embargo, los errores
  graves de cotizar muy bajo una propiedad de lujo están siendo
  controlados satisfactoriamente por el modelo. A pesar de su
  simplicidad e independencia condicional, Naive Bayes supera el azar
  estadístico, aunque para estrategias de alta precisión financiera
  requerirá apoyo de otros métodos.
\end{itemize}

\end{document}
